%! Minimizing Loss by Gradient Descent — Unit 8, Lesson 8.3 — Slides (Beamer)
\documentclass[aspectratio=169, 11pt]{beamer}
\usepackage{linalg-beamer}   % provides \burgundyheader{} and \sectionlabel[color]{}
\newcommand{\vv}{\mathbf{v}}
\newcommand{\bb}{\mathbf{b}}
\newcommand{\ww}{\mathbf{w}}
\newcommand{\relu}{\mathrm{ReLU}}

\begin{document}

% TITLE SLIDE (CourseName is not defined in beamer, so the name is literal)
{
\setbeamercolor{background canvas}{bg=burgundy}
\begin{frame}[plain]
  \vspace{2.8cm}\hspace{0.55cm}%
  \begin{minipage}{0.9\textwidth}
    {\color{goldacc}\bfseries\small LESSON 8.3}\\[0.5cm]
    {\color{white}\bfseries\LARGE Minimizing Loss by Gradient Descent}\\[1.0cm]
    {\color{blushmid}\small Linear Algebra~$\cdot$~Unit 8: Learning from Data}
  \end{minipage}
\end{frame}
}

% HOOK
\begin{frame}[t, plain]
  \burgundyheader{We can build networks --- but which weights?}
  \sectionlabel{Hook}
  \begin{columns}[T]
    \begin{column}{0.54\textwidth}
      We can \emph{build} networks: layers $\relu(A\vv+\bb)$ (8.0), folds (8.1), filters (8.2).\\[6pt]
      Every choice of weights gives a \textbf{loss} --- the Unit-4 sum of squared errors. Picture it as a
      \textbf{bowl} over the weights.\\[6pt]
      \textbf{Learning} $=$ find the \emph{bottom} (smallest loss). But there are \emph{millions} of weights.
    \end{column}
    \begin{column}{0.46\textwidth}
      \centering
      \begin{tikzpicture}[scale=0.5]
        \draw[->, charcoal, thin] (-0.3,0)--(5.2,0) node[right]{\scriptsize $w$};
        \draw[->, charcoal, thin] (0,-0.3)--(0,3.0) node[above]{\scriptsize $L(w)$};
        \draw[ultra thick, royalblue, smooth]
          plot coordinates {(0,2.4)(0.6,1.35)(1.2,0.6)(1.8,0.15)(2.4,0)(3.0,0.15)(3.6,0.6)(4.2,1.35)(4.8,2.4)};
        \fill[burgundy] (0.6,1.35) circle (2.4pt);
        \draw[->, very thick, goldacc] (1.05,1.15)--(2.15,0.2);
        \node[goldacc] at (2.55,0.9) {\scriptsize downhill};
        \fill[charcoal] (2.4,0) circle (1.6pt);
        \node[charcoal] at (2.4,-0.7) {\scriptsize best $w$};
      \end{tikzpicture}
    \end{column}
  \end{columns}
  \vspace{2pt}
  \begin{block}{The question of Lesson 8.3}
    How can a network tune millions of weights just by feeling which way is \textbf{downhill}?
  \end{block}
\end{frame}

% SLOPE = DOWNHILL
\begin{frame}[t, plain]
  \burgundyheader{Which way is downhill? The slope}
  \sectionlabel[goldacc]{Idea}
  \begin{columns}[T]
    \begin{column}{0.54\textwidth}
      The \textbf{slope} $L'(w)$ points \emph{uphill}. To lower the loss, step the \emph{opposite} way.\\[6pt]
      Our bowl $L(w)=(w-4)^2$ has slope
      \[ L'(w)=2w-8. \]
      At $w=0$: $L'(0)=-8$ (negative $\Rightarrow$ go right, toward $4$).\\[4pt]
      At the bottom $w=4$: $L'(4)=0$ --- flat ground, nowhere lower.
    \end{column}
    \begin{column}{0.46\textwidth}
      \centering
      \begin{tikzpicture}[scale=0.5]
        \draw[->, charcoal, thin] (-0.3,0)--(5.2,0) node[right]{\scriptsize $w$};
        \draw[->, charcoal, thin] (0,-0.3)--(0,3.0) node[above]{\scriptsize $L(w)$};
        \draw[ultra thick, royalblue, smooth]
          plot coordinates {(0,2.4)(0.6,1.35)(1.2,0.6)(1.8,0.15)(2.4,0)(3.0,0.15)(3.6,0.6)(4.2,1.35)(4.8,2.4)};
        \fill[burgundy] (0,2.4) circle (2.2pt);   \node[burgundy] at (-0.3,2.6) {\scriptsize $w_0$};
        \fill[burgundy] (1.2,0.6) circle (2.2pt); \node[burgundy] at (1.2,0.95) {\scriptsize $w_1$};
        \fill[burgundy] (1.8,0.15) circle (2.2pt);\node[burgundy] at (1.9,0.5) {\scriptsize $w_2$};
        \fill[burgundy] (2.1,0.04) circle (2.2pt);
        \draw[->, goldacc, thick] (0.25,2.1)--(1.05,0.8);
        \draw[->, goldacc, thick] (1.32,0.45)--(1.72,0.22);
        \node[charcoal] at (3.2,1.2) {\scriptsize each step lower};
      \end{tikzpicture}
    \end{column}
  \end{columns}
\end{frame}

% THE UPDATE RULE
\begin{frame}[t, plain]
  \burgundyheader{Gradient descent: step against the slope}
  \sectionlabel{Skill}
  \[ \boxed{\,w \;\leftarrow\; w - s\,L'(w)\,}\qquad s=\text{\textbf{learning rate} (step size)} \]
  Roll down $L(w)=(w-4)^2$ from $w=0$ with $s=0.25$:
  \begin{center}
  \renewcommand{\arraystretch}{1.25}
  \begin{tabular}{c|c|c}
  $w$ & slope $2w-8$ & new $w=w-0.25\,L'(w)$\\ \hline
  $0$ & $-8$ & $0-0.25(-8)=\mathbf{2}$\\
  $2$ & $-4$ & $2-0.25(-4)=\mathbf{3}$\\
  $3$ & $-2$ & $3-0.25(-2)=\mathbf{3.5}$\\
  \end{tabular}
  \end{center}
  \vspace{-2pt}
  The weights march $0\to2\to3\to3.5\to\cdots\to\mathbf{4}$; the loss falls $16\to4\to1\to0.25$.
  \textbf{Stop} when the slope is $0$ --- the step changes nothing.
\end{frame}

% MANY WEIGHTS + WHY IT MATTERS
\begin{frame}[t, plain]
  \burgundyheader{Many weights: the gradient --- and the whole picture}
  \sectionlabel[goldacc]{Payoff}
  \begin{itemize}
    \setlength{\itemsep}{6pt}
    \item The \textbf{gradient} $\nabla L$ is the slopes for \emph{all} weights at once. Same rule, a whole
          vector: $\ww \leftarrow \ww - s\,\nabla L$. \; E.g.\ $\nabla L=[2(w_1{-}1),2(w_2{-}3)]$ at $(0,0)$ is
          $[-2,-6]$; one step $\to(0.5,1.5)$, toward $(1,3)$.
    \item \textbf{Learning rate:} too small crawls; too big \emph{overshoots} the bottom (from $w=0$, $s=1$
          jumps to $w=8$).
    \item \textbf{The whole picture:} a network is Unit~1 (layers) $+$ one bend (ReLU) $+$ Unit~4 (loss) $+$
          rolling downhill (gradient descent) --- \emph{no new machinery}, repeated millions of times.
  \end{itemize}
  \vspace{2pt}
  \begin{block}{The road ahead}
    \textbf{8.4} mean, variance, and covariance --- the shape of the data itself.
  \end{block}
\end{frame}

\end{document}
